\chapter{Théorie des graphes}

\noindent\textit{Un graphe est un dessin de points reliés par des traits : il modélise
des \textbf{relations} entre des objets — villes reliées par des routes, joueurs s'affrontant
dans un tournoi, tâches d'un planning, examens incompatibles. Outil récent du programme de la
série C, la théorie des graphes relie l'\textbf{algèbre} (matrices) à des situations très
concrètes, et fournit des méthodes efficaces et élégantes pour étudier la connexité, les
parcours, le plus court chemin (algorithme de Dijkstra) et la coloration. Ce chapitre rassemble
tout le vocabulaire et tous les outils attendus au Baccalauréat.}
\vspace{8pt}

% ═══════════════════════════════════════════════════════════════
\section{Graphes, sommets et arêtes}

\subsection{Vocabulaire de base}

\begin{definition}[Graphe non orienté]
Un \textbf{graphe} (non orienté) $G=(S,A)$ est la donnée d'un ensemble fini $S$ de
\textbf{sommets} et d'un ensemble $A$ d'\textbf{arêtes}, chaque arête reliant deux sommets.
L'\textbf{ordre} du graphe est le nombre de sommets $\Card(S)$. Deux sommets reliés par une
arête sont dits \textbf{adjacents} (ou \textbf{voisins}) ; une arête reliant un sommet à
lui-même est une \textbf{boucle}.
\end{definition}

\begin{exemple}
Le graphe ci-dessous a pour sommets $S=\{A,B,C,D\}$ (ordre $4$) et pour arêtes
$\{AB,\,AC,\,BC,\,CD\}$ (quatre arêtes). Les sommets $A$ et $B$ sont adjacents ; $A$ et $D$
ne le sont pas.
\end{exemple}

\begin{illus}
\begin{tikzpicture}[>=Stealth,
  som/.style={circle,draw=onbo,fill=onbsoft,inner sep=2.5pt,font=\small\bfseries}]
\node[som] (A) at (0,2) {A};
\node[som] (B) at (2.4,2) {B};
\node[som] (C) at (1.2,0.4) {C};
\node[som] (D) at (3.4,-0.4) {D};
\draw[onbink,line width=0.9pt] (A) -- node[above,font=\scriptsize,onbmuted]{$e_1$} (B);
\draw[onbink,line width=0.9pt] (A) -- node[left,font=\scriptsize,onbmuted]{$e_2$} (C);
\draw[onbink,line width=0.9pt] (B) -- node[right,font=\scriptsize,onbmuted]{$e_3$} (C);
\draw[onbink,line width=0.9pt] (C) -- node[below,font=\scriptsize,onbmuted]{$e_4$} (D);
\node[align=left,text width=4.6cm,anchor=west] at (4.6,0.8)
{\small Ordre $4$, quatre arêtes. $A$ et $B$ sont \emph{adjacents} ; $C$ est voisin de
$A$, $B$ et $D$.};
\end{tikzpicture}
\fig{Figure 23.1 — Un graphe non orienté : sommets, arêtes étiquetées.}
\end{illus}

\subsection{Degré d'un sommet}

\begin{definition}[Degré]
Le \textbf{degré} d'un sommet $s$, noté $d(s)$, est le nombre d'arêtes qui ont $s$ pour
extrémité (une boucle compte pour $2$). Un sommet de degré $0$ est dit \textbf{isolé} ;
un sommet de degré $1$ est dit \textbf{pendant}.
\end{definition}

\begin{theoreme}[Lemme des poignées de main]
Dans tout graphe, la somme des degrés de tous les sommets est égale au double du nombre
d'arêtes :
\[
\sum_{s\in S} d(s) = 2\,\Card(A).
\]
\end{theoreme}

\noindent\emph{Justification.} Chaque arête possède exactement deux extrémités : elle est donc
comptée \emph{deux fois} dans la somme des degrés, une fois pour chacune. (Une boucle relie un
sommet à lui-même : elle apporte bien $2$ à son degré.)

\begin{propriete}[Nombre de sommets de degré impair]
Dans tout graphe, le nombre de sommets de degré \textbf{impair} est \textbf{pair}.
\end{propriete}

\begin{remarque}
La somme $\sum d(s)=2\,\Card(A)$ étant paire, les degrés impairs doivent se compenser :
ils sont donc en nombre pair. Conséquence pratique : dans un groupe, le nombre de personnes
ayant serré un nombre impair de mains est toujours pair.
\end{remarque}

\begin{exemple}
Sur la figure 23.1 : $d(A)=2$, $d(B)=2$, $d(C)=3$, $d(D)=1$. La somme vaut
$2+2+3+1=8=2\times 4$, ce qui confirme qu'il y a bien $4$ arêtes. Les sommets de degré
impair sont $C$ et $D$ : ils sont au nombre de $2$, qui est pair.
\end{exemple}

\begin{aretenir}
Pour tout graphe : $\displaystyle\sum_{s} d(s)=2\,\Card(A)$. Conséquences immédiates : la somme
des degrés est toujours \textbf{paire} ; le nombre de sommets de degré \textbf{impair} est
\textbf{pair} ; le nombre d'arêtes est la \emph{demi-somme} des degrés.
\end{aretenir}

% ═══════════════════════════════════════════════════════════════
\section{Graphes particuliers}

\begin{definition}[Graphe simple, complet, biparti, régulier]
\begin{itemize}
\item Un graphe est \textbf{simple} s'il est sans boucle et si deux sommets sont reliés
par \emph{au plus une} arête.
\item Un graphe simple est \textbf{complet} d'ordre $n$ (noté $K_n$) si \emph{toute} paire
de sommets est reliée par une arête.
\item Un graphe est \textbf{biparti} si l'on peut partager ses sommets en deux groupes
$X$ et $Y$ de sorte que toute arête joigne un sommet de $X$ à un sommet de $Y$ (aucune
arête à l'intérieur d'un groupe).
\item Un graphe est \textbf{$k$-régulier} si tous ses sommets ont le même degré $k$.
\end{itemize}
\end{definition}

\begin{propriete}[Nombre d'arêtes d'un graphe complet]
Le graphe complet $K_n$ possède
\[
\Card(A)=\binom{n}{2}=\frac{n(n-1)}{2}
\]
arêtes, et chacun de ses sommets est de degré $n-1$ : $K_n$ est $(n-1)$-régulier.
\end{propriete}

\begin{exemple}
$K_4$ (quatre sommets tous reliés) a $\dfrac{4\times 3}{2}=6$ arêtes, chaque sommet étant
de degré $3$. Un tournoi où $5$ équipes se rencontrent toutes une fois comporte
$\dfrac{5\times 4}{2}=10$ matches : c'est $K_5$.
\end{exemple}

\begin{illus}
\begin{tikzpicture}[>=Stealth,
  som/.style={circle,draw=onbo,fill=onbsoft,inner sep=2pt,font=\scriptsize\bfseries},
  somb/.style={circle,draw=onbblue,fill=white,inner sep=2pt,font=\scriptsize\bfseries}]
% K4
\begin{scope}
\foreach \i/\ang in {1/135,2/45,3/-45,4/-135}{
  \node[som] (k\i) at (\ang:1) {$\i$};}
\foreach \i/\j in {1/2,1/3,1/4,2/3,2/4,3/4}{
  \draw[onbink,line width=0.7pt] (k\i)--(k\j);}
\node[onbmuted,font=\small] at (0,-1.7) {$K_4$ : graphe complet};
\end{scope}
% biparti
\begin{scope}[xshift=6cm]
\node[som] (x1) at (0,1) {$x_1$};
\node[som] (x2) at (0,0) {$x_2$};
\node[som] (x3) at (0,-1) {$x_3$};
\node[somb] (y1) at (2,0.5) {$y_1$};
\node[somb] (y2) at (2,-0.5) {$y_2$};
\foreach \x in {x1,x2,x3}{\foreach \y in {y1,y2}{
  \draw[onbink,line width=0.6pt] (\x)--(\y);}}
\node[onbmuted,font=\small] at (1,-1.7) {biparti complet $K_{3,2}$};
\end{scope}
\end{tikzpicture}
\fig{Figure 23.2 — À gauche, le graphe complet $K_4$ ; à droite, un graphe biparti $K_{3,2}$.}
\end{illus}

\begin{remarque}
Un graphe est biparti si et seulement si tous ses cycles ont une longueur \textbf{paire}.
En particulier, un triangle (cycle de longueur $3$) ne peut jamais apparaître dans un graphe
biparti.
\end{remarque}

\subsection{Graphes orientés et pondérés}

\begin{definition}[Graphe orienté]
Un \textbf{graphe orienté} $G=(S,A)$ est un graphe dont les arêtes, appelées \textbf{arcs},
ont un \emph{sens} : un arc $(s_i,s_j)$ va de $s_i$ \emph{vers} $s_j$ (on le dessine par une
flèche). On distingue alors :
\begin{itemize}
\item le \textbf{degré sortant} $d^{+}(s)$ : nombre d'arcs partant de $s$ ;
\item le \textbf{degré entrant} $d^{-}(s)$ : nombre d'arcs arrivant en $s$.
\end{itemize}
\end{definition}

\begin{definition}[Graphe pondéré]
Un \textbf{graphe pondéré} est un graphe (orienté ou non) dont chaque arête (ou arc) porte
un nombre positif appelé \textbf{poids} (longueur, durée, coût, distance\ldots). Le poids
d'une chaîne est la \emph{somme} des poids des arêtes qui la composent.
\end{definition}

\begin{illus}
\begin{tikzpicture}[>=Stealth,
  som/.style={circle,draw=onbo,fill=onbsoft,inner sep=2.5pt,font=\small\bfseries}]
% graphe orienté
\begin{scope}
\node[som] (A) at (0,1.6) {A};
\node[som] (B) at (2,1.6) {B};
\node[som] (C) at (1,0) {C};
\draw[->,onbink,line width=0.9pt] (A)--(B);
\draw[->,onbink,line width=0.9pt] (B)--(C);
\draw[->,onbink,line width=0.9pt] (C)--(A);
\node[onbmuted,font=\small] at (1,-0.9) {orienté};
\end{scope}
% graphe pondéré
\begin{scope}[xshift=6cm]
\node[som] (P) at (0,1.6) {P};
\node[som] (Q) at (2.2,1.6) {Q};
\node[som] (R) at (1.1,0) {R};
\draw[onbink,line width=0.9pt] (P)-- node[above,font=\scriptsize,onbblue]{$5$} (Q);
\draw[onbink,line width=0.9pt] (P)-- node[left,font=\scriptsize,onbblue]{$3$} (R);
\draw[onbink,line width=0.9pt] (R)-- node[right,font=\scriptsize,onbblue]{$4$} (Q);
\node[onbmuted,font=\small] at (1.1,-0.9) {pondéré};
\end{scope}
\end{tikzpicture}
\fig{Figure 23.3 — À gauche un graphe orienté (arcs fléchés) ; à droite un graphe pondéré (poids sur les arêtes).}
\end{illus}

\begin{remarque}
Dans un graphe orienté, le lemme des poignées de main devient
$\sum_s d^{+}(s)=\sum_s d^{-}(s)=\Card(A)$ : chaque arc compte une fois comme sortant et une
fois comme entrant. Les graphes pondérés servent à modéliser les \textbf{réseaux} (routes avec
distances, tâches avec durées) ; ils sont au cœur du problème du \emph{plus court chemin}.
\end{remarque}

% ═══════════════════════════════════════════════════════════════
\section{Matrice d'adjacence}

\begin{definition}[Matrice d'adjacence]
On numérote les sommets $s_1,s_2,\dots,s_n$. La \textbf{matrice d'adjacence} de $G$ est la
matrice carrée $M=(m_{ij})$ d'ordre $n$ définie par
\[
m_{ij}=\text{nombre d'arêtes (ou d'arcs) reliant } s_i \text{ à } s_j.
\]
Pour un graphe simple, $m_{ij}\in\{0,1\}$ et $m_{ii}=0$ (pas de boucle).
\end{definition}

\begin{propriete}[Lecture de la matrice]
Pour un graphe \textbf{non orienté}, la matrice d'adjacence est \textbf{symétrique}
($m_{ij}=m_{ji}$) ; la somme des coefficients de la ligne $i$ (ou de la colonne $i$) est
égale au \textbf{degré} du sommet $s_i$. Pour un graphe \textbf{orienté}, $M$ n'est en
général pas symétrique : la somme de la ligne $i$ donne le degré \emph{sortant} $d^{+}(s_i)$,
celle de la colonne $i$ le degré \emph{entrant} $d^{-}(s_i)$.
\end{propriete}

\begin{illus}
\begin{tikzpicture}[>=Stealth,
  som/.style={circle,draw=onbo,fill=onbsoft,inner sep=2.5pt,font=\small\bfseries}]
\node[som] (A) at (0,2) {A};
\node[som] (B) at (2,2) {B};
\node[som] (C) at (0,0) {C};
\node[som] (D) at (2,0) {D};
\draw[onbink,line width=0.9pt] (A)--(B);
\draw[onbink,line width=0.9pt] (A)--(C);
\draw[onbink,line width=0.9pt] (B)--(C);
\draw[onbink,line width=0.9pt] (B)--(D);
\node[anchor=west] at (4,1)
{$\displaystyle
M=\begin{array}{c@{\,}c}
 & \begin{array}{cccc}\!A\!&\!B\!&\!C\!&\!D\!\end{array}\\[2pt]
\begin{array}{c}A\\B\\C\\D\end{array} &
\left(\begin{array}{cccc}
0&1&1&0\\
1&0&1&1\\
1&1&0&0\\
0&1&0&0
\end{array}\right)
\end{array}$};
\end{tikzpicture}
\fig{Figure 23.4 — Un graphe et sa matrice d'adjacence (ordre des sommets $A,B,C,D$).}
\end{illus}

\begin{exemple}
Sur la figure 23.4 : la ligne de $B$ est $(1,0,1,1)$, de somme $3=d(B)$. La matrice est bien
symétrique. La diagonale est nulle car le graphe est simple.
\end{exemple}

\subsection{Puissances de la matrice d'adjacence}

\begin{theoreme}[Nombre de chaînes de longueur $k$]
Soit $M$ la matrice d'adjacence d'un graphe. Pour tout entier $k\ge 1$, le coefficient
$(i,j)$ de la matrice $M^k$ est égal au \textbf{nombre de chaînes de longueur $k$} reliant
le sommet $s_i$ au sommet $s_j$.
\end{theoreme}

\begin{aretenir}
Les \textbf{puissances de $M$} comptent les chemins : $M^2$ donne le nombre de trajets en
\emph{deux} arêtes, $M^3$ en \emph{trois}, etc. C'est l'outil-clé pour répondre à
« de combien de façons peut-on aller de $X$ à $Y$ en exactement $k$ étapes ? ». Sur la
diagonale, $(M^2)_{ii}=d(s_i)$ : le nombre de chaînes de longueur $2$ d'un sommet à lui-même
est son nombre de voisins.
\end{aretenir}

\begin{exemple}
Avec la matrice $M$ de la figure 23.4, on calcule
\[
M^2=\begin{pmatrix}2&1&1&1\\1&3&1&0\\1&1&2&1\\1&0&1&1\end{pmatrix}.
\]
Le coefficient $(A,D)$ vaut $1$ : il existe \emph{une} seule chaîne de longueur $2$ de $A$
à $D$, à savoir $A-B-D$. Le coefficient $(B,B)$ vaut $3$ : $B$ a trois voisins ($A,C,D$),
donc trois allers-retours de longueur $2$.
\end{exemple}

% ═══════════════════════════════════════════════════════════════
\section{Chaînes, cycles et connexité}

\begin{definition}[Chaîne, cycle]
Une \textbf{chaîne} est une suite de sommets reliés de proche en proche par des arêtes ;
sa \textbf{longueur} est le nombre d'arêtes parcourues. Une chaîne est \textbf{simple} si
elle n'emprunte pas deux fois la même arête. Un \textbf{cycle} est une chaîne simple
\emph{fermée} (même sommet de départ et d'arrivée) de longueur au moins $3$. Dans un graphe
orienté, on parle de \textbf{chemin} et de \textbf{circuit}.
\end{definition}

\begin{definition}[Graphe connexe]
Un graphe est \textbf{connexe} si, pour tout couple de sommets, il existe une chaîne les
reliant : « on peut aller de n'importe quel sommet à n'importe quel autre ». Sinon, les
sommets se répartissent en \textbf{composantes connexes} (les « morceaux » du graphe).
\end{definition}

\begin{propriete}[Connexité par la matrice]
Un graphe d'ordre $n$ est connexe si et seulement si la matrice
\[
B=I_n+M+M^2+\dots+M^{n-1}
\]
ne comporte \textbf{aucun coefficient nul}. ($I_n$ tient compte du « trajet vide » d'un
sommet à lui-même.)
\end{propriete}

\begin{remarque}
Un coefficient $b_{ij}$ de $B$ compte le nombre total de chaînes de longueur $\le n-1$ de
$s_i$ à $s_j$. Si $s_i$ et $s_j$ sont reliés, ils le sont par une chaîne d'au plus $n-1$
arêtes (sans repasser par un sommet) : $b_{ij}\ne0$. C'est pourquoi tester la connexité
revient à vérifier que $B$ n'a aucun zéro.
\end{remarque}

\subsection{Chaînes et cycles eulériens}

\begin{definition}[Chaîne / cycle eulérien]
Une \textbf{chaîne eulérienne} est une chaîne qui emprunte \emph{une et une seule fois}
\textbf{chaque arête} du graphe. Si elle est fermée (retour au point de départ), on parle
de \textbf{cycle eulérien}.
\end{definition}

\begin{theoreme}[Théorème d'Euler]
Soit $G$ un graphe connexe.
\begin{itemize}
\item $G$ admet un \textbf{cycle eulérien} si et seulement si \emph{tous} ses sommets sont
de degré \textbf{pair}.
\item $G$ admet une \textbf{chaîne eulérienne} (non fermée) si et seulement s'il a
\emph{exactement deux} sommets de degré impair ; le parcours commence alors à l'un de ces
deux sommets et finit à l'autre.
\end{itemize}
\end{theoreme}

\begin{illus}
\begin{tikzpicture}[>=Stealth,
  som/.style={circle,draw=onbo,fill=onbsoft,inner sep=2.5pt,font=\small\bfseries}]
% les ponts de Königsberg, schématisés (4 sommets, multi-arêtes)
\node[som] (A) at (0,1) {A};
\node[som] (B) at (3,1) {B};
\node[som] (C) at (1.5,2.4) {C};
\node[som] (D) at (1.5,-0.4) {D};
\draw[onbink,line width=0.8pt] (A) to[bend left=25] (C);
\draw[onbink,line width=0.8pt] (A) to[bend right=25] (C);
\draw[onbink,line width=0.8pt] (B) to[bend left=25] (C);
\draw[onbink,line width=0.8pt] (B) to[bend right=25] (C);
\draw[onbink,line width=0.8pt] (A) -- (D);
\draw[onbink,line width=0.8pt] (B) -- (D);
\draw[onbink,line width=0.8pt] (C) -- (D);
\node[align=left,text width=4.6cm,anchor=west] at (4.2,1)
{\small Königsberg : $d(C)=5$, $d(A)=d(B)=d(D)=3$. Quatre sommets de degré \emph{impair}
$\Rightarrow$ aucune chaîne eulérienne : la promenade est \textbf{impossible}.};
\end{tikzpicture}
\fig{Figure 23.5 — Les sept ponts de Königsberg modélisés par un (multi)graphe.}
\end{illus}

\begin{exemple}
Le problème historique des \emph{ponts de Königsberg} (figure 23.5) revient à chercher une
chaîne eulérienne dans un graphe ayant quatre sommets tous de degré impair : c'est donc
\textbf{impossible}. Un facteur peut parcourir toutes les rues d'un quartier sans repasser
deux fois par la même seulement si au plus deux carrefours sont de degré impair.
\end{exemple}

\subsection{Plus court chemin : algorithme de Dijkstra}

\begin{definition}[Plus court chemin]
Dans un graphe \textbf{pondéré} à poids positifs, le \textbf{plus court chemin} entre deux
sommets est une chaîne reliant ces sommets dont le poids total (somme des poids des arêtes)
est \emph{minimal}.
\end{definition}

\begin{theoreme}[Algorithme de Dijkstra]
Pour trouver le plus court chemin d'un sommet de départ $D$ vers tous les autres dans un
graphe à poids positifs :
\begin{enumerate}[label=\arabic*)]
\item Affecter au départ $D$ la distance $0$ et à tous les autres sommets la distance $+\infty$
(« provisoire »).
\item Choisir parmi les sommets \emph{non encore figés} celui de plus petite distance provisoire ;
le \textbf{figer} (sa distance est définitive).
\item \textbf{Mettre à jour} chaque voisin $V$ non figé : si \,(distance du sommet figé)\,$+$\,(poids
de l'arête)\, est inférieure à la distance provisoire de $V$, la remplacer, et noter le
\emph{prédécesseur} de $V$.
\item Répéter tant qu'il reste des sommets non figés.
\end{enumerate}
La distance finale d'un sommet est le poids du plus court chemin depuis $D$ ; on le reconstitue
en remontant les prédécesseurs.
\end{theoreme}

\begin{illus}
\begin{tikzpicture}[>=Stealth,
  som/.style={circle,draw=onbo,fill=onbsoft,inner sep=2.5pt,font=\small\bfseries}]
\node[som] (A) at (0,1.5) {A};
\node[som] (B) at (2.2,2.6) {B};
\node[som] (C) at (2.2,0.4) {C};
\node[som] (D) at (4.4,1.5) {D};
\node[som] (E) at (6.3,1.5) {E};
\draw[onbink,line width=0.8pt] (A)-- node[above left,font=\scriptsize,onbblue]{$4$} (B);
\draw[onbink,line width=0.8pt] (A)-- node[below left,font=\scriptsize,onbblue]{$2$} (C);
\draw[onbink,line width=0.8pt] (B)-- node[left,font=\scriptsize,onbblue]{$1$} (C);
\draw[onbink,line width=0.8pt] (B)-- node[above,font=\scriptsize,onbblue]{$5$} (D);
\draw[onbink,line width=0.8pt] (C)-- node[below,font=\scriptsize,onbblue]{$8$} (D);
\draw[onbink,line width=0.8pt] (D)-- node[above,font=\scriptsize,onbblue]{$3$} (E);
\node[onbmuted,font=\scriptsize] at (3.1,-0.6)
{Plus court chemin $A\to E$ : $A-C-B-D-E$ de poids $2+1+5+3=11$.};
\end{tikzpicture}
\fig{Figure 23.6 — Recherche d'un plus court chemin par l'algorithme de Dijkstra.}
\end{illus}

\begin{exemple}
Sur la figure 23.6, cherchons le plus court chemin de $A$ à $E$. On organise le calcul en
tableau (distance provisoire la plus faible \underline{soulignée} = sommet figé) :
\[
\begin{array}{c|ccccc|l}
\text{étape} & A & B & C & D & E & \text{figé}\\\hline
0 & \underline{0} & \infty & \infty & \infty & \infty & A\\
1 & & 4 & \underline{2}_{(A)} & \infty & \infty & C\\
2 & & \underline{3}_{(C)} & & 10 & \infty & B\\
3 & & & & \underline{8}_{(B)} & \infty & D\\
4 & & & & & \underline{11}_{(D)} & E
\end{array}
\]
On lit en remontant les prédécesseurs : $E\leftarrow D\leftarrow B\leftarrow C\leftarrow A$,
soit le chemin $A-C-B-D-E$ de longueur totale $\mathbf{11}$.
\end{exemple}

\subsection{Coloration d'un graphe}

\begin{definition}[Coloration, nombre chromatique]
\textbf{Colorer} un graphe, c'est attribuer une couleur à chaque sommet de sorte que deux
sommets \textbf{adjacents} reçoivent des couleurs \emph{différentes}. Le \textbf{nombre
chromatique} $\chi(G)$ est le plus petit nombre de couleurs permettant une telle coloration.
\end{definition}

\begin{propriete}[Encadrement du nombre chromatique]
Si $\Delta$ est le degré maximal des sommets, alors $\chi(G)\le \Delta+1$. De plus, si le
graphe contient un sous-graphe complet $K_p$ (une « clique » de $p$ sommets tous reliés),
alors $\chi(G)\ge p$. En particulier $\chi(K_n)=n$ et $\chi(G)=2$ pour tout graphe biparti
(possédant au moins une arête).
\end{propriete}

\begin{remarque}
La coloration modélise les problèmes de \textbf{compatibilité} : couleurs = créneaux
horaires, fréquences radio, ou groupes ; une arête = « ces deux objets ne peuvent pas
partager le même créneau ». Le nombre chromatique donne alors le nombre minimal de créneaux.
L'\textbf{algorithme de Welsh-Powell} (colorer les sommets par degré décroissant, avec la
plus petite couleur libre) fournit une coloration valable, souvent optimale.
\end{remarque}

\begin{aretenir}
Pour le nombre chromatique : on \textbf{minore} $\chi$ par la taille de la plus grosse clique
(un triangle force $\chi\ge3$), et on \textbf{majore} $\chi$ en exhibant une coloration
explicite. Quand la minoration et la majoration coïncident, on a la valeur exacte de $\chi$.
\end{aretenir}

% ═══════════════════════════════════════════════════════════════
\section{L'essentiel à retenir}

\begin{synthese}
\textbf{Vocabulaire.} Graphe $G=(S,A)$ : ordre $=\Card(S)$, $\Card(A)$ arêtes. Degré
$d(s)=$ nombre d'arêtes incidentes (boucle $=2$). Orienté : degrés $d^{+}$ (sortant),
$d^{-}$ (entrant). Pondéré : poids sur les arêtes.\\[4pt]
\textbf{Poignées de main.} $\displaystyle\sum_s d(s)=2\,\Card(A)$ : la somme des degrés est
paire ; le nombre de sommets de degré impair est pair.\\[4pt]
\textbf{Graphes types.} $K_n$ complet : $\tfrac{n(n-1)}2$ arêtes, $(n-1)$-régulier,
$\chi=n$. Biparti : pas de cycle impair, $\chi=2$.\\[4pt]
\textbf{Matrice $M$.} $m_{ij}=$ nb d'arêtes $s_i\!-\!s_j$ ; symétrique si non orienté ; somme
d'une ligne $=$ degré. $(M^k)_{ij}=$ nombre de chaînes de longueur $k$ de $s_i$ à $s_j$ ;
$(M^2)_{ii}=d(s_i)$.\\[4pt]
\textbf{Connexité.} On peut relier tout couple de sommets. Test : $B=I_n+M+\dots+M^{n-1}$ sans
aucun zéro $\iff$ connexe.\\[4pt]
\textbf{Euler (graphe connexe).} \emph{cycle} eulérien $\iff$ tous les degrés pairs ;
\emph{chaîne} eulérienne $\iff$ exactement deux degrés impairs (départ/arrivée à ces deux
sommets) ; sinon, aucun parcours eulérien.\\[4pt]
\textbf{Dijkstra (poids $\ge0$).} départ à $0$, autres à $+\infty$ ; figer le plus petit ;
mettre à jour les voisins (dist.\ figée $+$ poids), noter les prédécesseurs ; répéter ;
remonter pour le chemin.\\[4pt]
\textbf{Coloration.} $\chi(G)=$ nb minimal de couleurs (sommets adjacents $\ne$). Clique $K_p$
$\Rightarrow\chi\ge p$ ; $\chi\le\Delta+1$. Applications : emplois du temps, fréquences.\\[4pt]
\textbf{Pièges.} oublier qu'une boucle vaut $2$ au degré ; lire un degré sur la colonne au lieu
de la ligne ; confondre cycle (degrés pairs) et chaîne (deux impairs) eulériens ; appliquer
Dijkstra avec des poids négatifs (interdit).
\end{synthese}

% ═══════════════════════════════════════════════════════════════
\section{Méthodes \& savoir-faire}

\methode{Modéliser une situation concrète par un graphe}
Choisir ce que représentent les \textbf{sommets} (objets) et les \textbf{arêtes} (la relation
qui les lie), puis dénombrer ordre et degrés.
\begin{exemple}
Cinq villes — Yaoundé, Douala, Bafoussam, Bertoua, Garoua — et les routes directes :
Yaoundé–Douala, Yaoundé–Bafoussam, Yaoundé–Bertoua, Douala–Bafoussam, Bafoussam–Garoua.
On pose un sommet par ville, une arête par route directe : graphe d'ordre $5$, $5$ arêtes.
Degrés : $d(\text{Yaoundé})=3$, $d(\text{Bafoussam})=3$, $d(\text{Douala})=2$,
$d(\text{Garoua})=1$, $d(\text{Bertoua})=1$. Vérification : $3+3+2+1+1=10=2\times 5$.
\end{exemple}

\methode{Utiliser le lemme des poignées de main}
La somme des degrés vaut $2\Card(A)$ : elle est paire, et on en déduit le nombre d'arêtes ou
l'impossibilité d'une configuration.
\begin{exemple}
Existe-t-il un graphe à $5$ sommets de degrés $3,3,3,3,3$ ? La somme $5\times3=15$ est
\emph{impaire} : impossible (elle devrait être $2\Card(A)$, donc paire). En revanche, des
degrés $3,3,2,2,2$ donnent une somme $12=2\times6$ : $6$ arêtes — c'est envisageable.
\end{exemple}

\methode{Écrire ou lire une matrice d'adjacence}
Fixer un ordre des sommets, puis remplir $m_{ij}=1$ si $s_i$ et $s_j$ sont reliés, $0$ sinon.
La matrice doit être symétrique (non orienté) ; la somme d'une ligne redonne le degré.
\begin{exemple}
Pour le triangle $A,B,C$ tous reliés deux à deux ($K_3$) :
$M=\begin{pmatrix}0&1&1\\1&0&1\\1&1&0\end{pmatrix}$. Chaque ligne somme à $2$ : chaque
sommet est de degré $2$, ce qui est correct pour $K_3$.
\end{exemple}

\methode{Compter les chaînes par les puissances de la matrice}
Calculer $M^k$ ; le coefficient $(i,j)$ donne le nombre de chaînes de longueur $k$ de $s_i$
à $s_j$.
\begin{exemple}
Pour $K_3$ ci-dessus, $M^2=\begin{pmatrix}2&1&1\\1&2&1\\1&1&2\end{pmatrix}$. Le coefficient
$(A,B)$ vaut $1$ : une seule chaîne de longueur $2$ de $A$ à $B$, à savoir $A-C-B$. Sur la
diagonale, $2=d(A)$.
\end{exemple}

\methode{Étudier la connexité et chercher une chaîne}
Partir d'un sommet et explorer de proche en proche tous les sommets atteignables ; si on
les atteint tous, le graphe est connexe. Sinon, on a identifié une \textbf{composante
connexe}.
\begin{exemple}
Graphe de sommets $\{1,2,3,4,5\}$ et arêtes $12,\,23,\,45$. En partant de $1$ on atteint
$1,2,3$ mais ni $4$ ni $5$ : le graphe \emph{n'est pas} connexe ; il a deux composantes
$\{1,2,3\}$ et $\{4,5\}$.
\end{exemple}

\methode{Appliquer le théorème d'Euler}
Compter les sommets de degré impair sur un graphe connexe : $0\Rightarrow$ cycle eulérien ;
exactement $2\Rightarrow$ chaîne eulérienne (entre ces deux sommets) ; sinon, aucun parcours
eulérien.
\begin{exemple}
Un graphe connexe de degrés $(2,2,4,2)$ : tous pairs $\Rightarrow$ cycle eulérien possible.
Degrés $(3,2,2,3)$ : exactement deux impairs $\Rightarrow$ chaîne eulérienne possible, qui
commence et finit aux deux sommets de degré $3$.
\end{exemple}

\methode{Appliquer l'algorithme de Dijkstra}
Dresser le tableau des distances : départ à $0$, autres à $+\infty$ ; à chaque étape, figer
le plus petit, mettre à jour les voisins (distance figée $+$ poids), noter les prédécesseurs.
\begin{exemple}
Graphe : $A\!-\!B$ (poids $1$), $A\!-\!C$ (poids $4$), $B\!-\!C$ (poids $1$). De $A$ :
on fige $A$ ($0$), puis $B$ ($1$). Par $B$, on atteint $C$ à $1+1=2<4$ : on fige $C$ à $2$,
prédécesseur $B$. Plus court chemin $A\to C$ : $A-B-C$, longueur $2$ (et non $4$ en direct).
\end{exemple}

\methode{Colorer un graphe et trouver son nombre chromatique}
Colorer les sommets un à un (en priorité ceux de plus haut degré, méthode de Welsh-Powell)
avec la plus petite couleur disponible ; minorer $\chi$ par la taille du plus grand
sous-graphe complet.
\begin{exemple}
Un graphe contenant un triangle a $\chi\ge 3$. S'il est de plus colorable avec $3$ couleurs,
alors $\chi=3$. Un graphe biparti (sans triangle, deux groupes) a $\chi=2$.
\end{exemple}

\methode{Résoudre un problème d'emploi du temps par coloration}
Modéliser : sommets $=$ épreuves/réunions, arête $=$ incompatibilité (ne peuvent pas être
simultanées). Le nombre chromatique $=$ nombre minimal de créneaux.
\begin{exemple}
Trois clubs se réunissent ; deux clubs partageant un membre ne peuvent se réunir le même
jour. Si les incompatibilités forment un triangle, il faut $\chi=3$ jours ; si elles forment
une simple chaîne, $\chi=2$ jours suffisent.
\end{exemple}

% ═══════════════════════════════════════════════════════════════
\section{Exercices d'application}

\rubrique{Sommets, arêtes et degrés}
\exo{}\diff{1} Un graphe a $6$ sommets de degrés respectifs $3,3,2,2,1,1$. Combien
possède-t-il d'arêtes ? Le nombre de sommets de degré impair est-il pair ?

\exo{}\diff{1} Dessiner, si possible, un graphe simple d'ordre $4$ dont tous les sommets
sont de degré $3$. Combien a-t-il d'arêtes ?

\exo{}\diff{2} Existe-t-il un graphe simple à $5$ sommets dont les degrés sont
$4,4,4,4,1$ ? Justifier.

\exo{}\diff{2} Dans une classe de Terminale C, on relie deux élèves dès qu'ils sont amis.
Peut-il exister exactement $7$ élèves ayant chacun un nombre \emph{impair} d'amis ? Justifier.

\rubrique{Graphes particuliers}
\exo{}\diff{1} Combien d'arêtes possède le graphe complet $K_7$ ? Quel est le degré de
chacun de ses sommets ?

\exo{}\diff{2} Dans un tournoi, $8$ équipes se rencontrent toutes exactement une fois.
Combien de matches sont joués au total ?

\exo{}\diff{2} Un graphe orienté a pour arcs $A\!\to\!B$, $B\!\to\!C$, $A\!\to\!C$,
$C\!\to\!A$. Donner les degrés entrant et sortant de chaque sommet et vérifier que
$\sum d^{+}=\sum d^{-}$.

\rubrique{Matrice d'adjacence}
\exo{}\diff{2} On donne la matrice $M$ d'un graphe de sommets $A,B,C,D$ :
\[
M=\begin{pmatrix}0&1&1&1\\1&0&0&1\\1&0&0&1\\1&1&1&0\end{pmatrix}.
\]
Dessiner le graphe, donner le degré de chaque sommet et dire s'il est connexe.

\exo{}\diff{3} Soit $M$ la matrice d'adjacence d'un graphe d'ordre $4$ (sommets
$A,B,C,D$). On calcule
\[
M^2=\begin{pmatrix}3&1&1&2\\1&2&2&1\\1&2&2&1\\2&1&1&3\end{pmatrix}.
\]
Donner le degré de chaque sommet et le nombre de chaînes de longueur $2$ entre $A$ et $D$.

\exo{}\diff{2} Pour le graphe de la figure 23.1 ($A\!-\!B$, $A\!-\!C$, $B\!-\!C$, $C\!-\!D$),
écrire la matrice d'adjacence $M$ dans l'ordre $A,B,C,D$ et vérifier qu'elle est symétrique.

\rubrique{Connexité et parcours eulériens}
\exo{}\diff{2} Un graphe connexe a pour degrés $5,5,4,2,2$. Admet-il un cycle eulérien ?
une chaîne eulérienne ? Préciser dans ce cas les sommets de départ et d'arrivée.

\exo{}\diff{2} Un dessin se compose d'une enveloppe : un rectangle $ABCD$ surmonté d'un
triangle (sommet $E$ au-dessus de $AB$), avec les diagonales $AC$ et $BD$. Peut-on le tracer
« sans lever le crayon ni repasser sur un trait » ? Justifier par le théorème d'Euler.

\exo{}\diff{3} Un facteur doit parcourir toutes les rues d'un quartier modélisé par un graphe
connexe dont les degrés sont $4,4,3,3,2,2$. Peut-il faire sa tournée en passant une seule fois
par chaque rue ? Si oui, doit-il revenir à son point de départ ?

\rubrique{Plus court chemin (Dijkstra)}
\exo{}\diff{2} Dans un graphe pondéré, on a les arêtes (et poids) : $A\!-\!B\,(2)$,
$A\!-\!C\,(5)$, $B\!-\!C\,(1)$, $B\!-\!D\,(7)$, $C\!-\!D\,(3)$. Déterminer, par l'algorithme
de Dijkstra, le plus court chemin de $A$ à $D$ et sa longueur.

\exo{}\diff{3} Quatre villes $P,Q,R,S$ sont reliées par des routes (en km) :
$P\!-\!Q\,(4)$, $P\!-\!R\,(2)$, $Q\!-\!R\,(1)$, $Q\!-\!S\,(5)$, $R\!-\!S\,(8)$. Déterminer
l'itinéraire le plus court de $P$ à $S$ et sa distance.

\rubrique{Coloration}
\exo{}\diff{2} On veut programmer $5$ examens $E_1,\dots,E_5$. Deux examens ne peuvent avoir
lieu en même temps si un même élève les passe tous deux. Les incompatibilités sont :
$E_1$–$E_2$, $E_1$–$E_3$, $E_2$–$E_3$, $E_3$–$E_4$, $E_4$–$E_5$. Modéliser par un graphe et
déterminer le nombre minimal de créneaux horaires.

\exo{}\diff{2} Déterminer le nombre chromatique du graphe complet $K_5$ et celui d'un cycle
de longueur $6$ (hexagone). Justifier.

\rubrique{Problème de synthèse}
\exo{}\diff{3} \emph{Réseau routier camerounais.} On considère cinq villes :
Yaoundé (Y), Douala (D), Bafoussam (B), Bertoua (R) et Ngaoundéré (N). Les liaisons
routières directes sont : Y–D, Y–B, Y–R, D–B, B–N, R–N.
\begin{enumerate}[label=\arabic*)]
\item Représenter cette situation par un graphe et donner le degré de chaque ville.
\item Écrire la matrice d'adjacence $M$ (ordre $Y,D,B,R,N$). Le graphe est-il connexe ?
\item Un transporteur souhaite emprunter chaque route exactement une fois. Est-ce possible ?
Si oui, entre quelles villes doit-il commencer et finir ?
\item On admet que le coefficient $(Y,N)$ de $M^2$ vaut $2$. Interpréter ce résultat
en termes d'itinéraire.
\end{enumerate}

% ═══════════════════════════════════════════════════════════════
\section{Exercices d'entraînement}

\rubrique{Vocabulaire, degrés, poignées de main}
\exo{}\diff{1} Un graphe a $5$ sommets de degrés $4,3,3,2,2$. Combien a-t-il d'arêtes ?

\exo{}\diff{1} Dessiner un graphe d'ordre $4$ ayant exactement $3$ arêtes et un sommet isolé.

\exo{}\diff{2} Existe-t-il un graphe simple d'ordre $6$ dont les degrés sont $5,5,5,5,5,5$ ?
Si oui, lequel ?

\exo{}\diff{2} Montrer que dans tout graphe simple d'ordre $n\ge2$, il existe au moins deux
sommets de même degré.

\exo{}\diff{2} Un graphe simple d'ordre $7$ a tous ses sommets de degré $3$. Est-ce possible ?
Justifier.

\exo{}\diff{1} Combien d'arêtes a un graphe $4$-régulier d'ordre $10$ ?

\rubrique{Graphes complets, bipartis, orientés}
\exo{}\diff{1} Combien d'arêtes possède $K_{10}$ ?

\exo{}\diff{2} Dans une réunion, $12$ personnes se serrent toutes la main une fois. Combien
de poignées de main au total ?

\exo{}\diff{2} Le graphe biparti complet $K_{3,4}$ : combien de sommets, combien d'arêtes ?
Quel est le degré de chaque sommet de chaque groupe ?

\exo{}\diff{2} Dessiner un graphe orienté d'ordre $4$ tel que $d^{+}(A)=2$ et $d^{-}(A)=1$.

\exo{}\diff{3} Montrer qu'un graphe biparti ne contient aucun triangle.

\rubrique{Matrices d'adjacence et puissances}
\exo{}\diff{1} Écrire la matrice d'adjacence du cycle $A\!-\!B\!-\!C\!-\!D\!-\!A$.

\exo{}\diff{2} On donne $M=\begin{pmatrix}0&1&1&0\\1&0&1&1\\1&1&0&1\\0&1&1&0\end{pmatrix}$
(sommets $A,B,C,D$). Dessiner le graphe et donner les degrés.

\exo{}\diff{2} Pour $K_3$ de matrice $M=\begin{pmatrix}0&1&1\\1&0&1\\1&1&0\end{pmatrix}$,
calculer $M^2$ et interpréter le coefficient $(A,B)$.

\exo{}\diff{3} Soit $M$ la matrice du graphe : $A\!-\!B$, $B\!-\!C$, $C\!-\!D$, $D\!-\!A$
(carré). Calculer $M^2$ et indiquer le nombre de chaînes de longueur $2$ de $A$ à $C$.

\exo{}\diff{3} Un graphe d'ordre $4$ a pour $M^2$ une diagonale $(2,3,3,2)$. Donner le degré
de chaque sommet et le nombre total d'arêtes.

\rubrique{Connexité et composantes}
\exo{}\diff{1} Le graphe de sommets $\{1,2,3,4,5,6\}$ et arêtes $12,23,45,56$ est-il connexe ?
Donner ses composantes connexes.

\exo{}\diff{2} Un graphe d'ordre $5$ a $3$ arêtes. Peut-il être connexe ? (On rappelle qu'un
graphe connexe d'ordre $n$ a au moins $n-1$ arêtes.)

\exo{}\diff{2} Combien de composantes connexes a un graphe d'ordre $8$ formé de deux triangles
disjoints et de deux sommets isolés ?

\exo{}\diff{3} Montrer qu'un graphe d'ordre $n$ ayant strictement plus de $\dfrac{(n-1)(n-2)}2$
arêtes est nécessairement connexe.

\rubrique{Parcours eulériens}
\exo{}\diff{1} Un graphe connexe a tous ses sommets de degré $4$. Admet-il un cycle eulérien ?

\exo{}\diff{2} Les degrés d'un graphe connexe sont $3,3,3,3$. Peut-on le tracer sans lever le
crayon ? Justifier.

\exo{}\diff{2} Un quartier a $6$ carrefours de degrés $4,4,2,2,3,3$. Le balayeur peut-il
parcourir chaque rue une seule fois ? Reviendra-t-il à son point de départ ?

\exo{}\diff{3} Pour quelles valeurs de $n$ le graphe complet $K_n$ admet-il un cycle eulérien ?
(Indication : étudier la parité de $n-1$.)

\rubrique{Plus court chemin}
\exo{}\diff{2} Arêtes pondérées : $A\!-\!B\,(3)$, $A\!-\!C\,(1)$, $C\!-\!B\,(1)$,
$B\!-\!D\,(2)$, $C\!-\!D\,(6)$. Plus court chemin de $A$ à $D$ ?

\exo{}\diff{3} Cinq villes $A,B,C,D,E$ et distances (km) : $A\!-\!B\,(6)$, $A\!-\!C\,(2)$,
$C\!-\!B\,(3)$, $B\!-\!E\,(1)$, $C\!-\!D\,(5)$, $D\!-\!E\,(2)$. Déterminer le plus court trajet
de $A$ à $E$.

\exo{}\diff{2} Dans le graphe pondéré $P\!-\!Q\,(4)$, $Q\!-\!R\,(2)$, $P\!-\!R\,(7)$,
appliquer Dijkstra de $P$ : quelles sont les distances minimales à $Q$ et à $R$ ?

\rubrique{Coloration}
\exo{}\diff{1} Quel est le nombre chromatique d'un cycle de longueur $5$ ?

\exo{}\diff{2} Six matières doivent être affectées à des salles ; deux matières partageant un
professeur ne peuvent pas avoir lieu en même temps. Les conflits sont $M_1M_2$, $M_2M_3$,
$M_3M_1$, $M_4M_5$, $M_5M_6$. Combien de créneaux au minimum ?

\exo{}\diff{2} Déterminer $\chi(G)$ pour un graphe formé d'un carré ($A\!-\!B\!-\!C\!-\!D\!-\!A$)
auquel on ajoute la diagonale $A\!-\!C$.

\exo{}\diff{3} On dispose des fréquences radio à attribuer à $7$ émetteurs ; deux émetteurs
trop proches (donc reliés par une arête) doivent recevoir des fréquences différentes. Le degré
maximal du graphe est $3$. Combien de fréquences suffisent à coup sûr ? Justifier par
l'encadrement $\chi\le\Delta+1$.

\rubrique{Problèmes de synthèse}
\exo{}\diff{3} On modélise un réseau de $5$ ordinateurs par un graphe d'arêtes $12,13,23,34,45$.
Donner sa matrice d'adjacence, ses degrés, dire s'il est connexe, et s'il admet une chaîne
eulérienne.

\exo{}\diff{3} Sept élèves doivent être répartis en équipes ; certaines paires refusent de
travailler ensemble (arêtes du graphe de conflit). Si ce graphe contient un $K_4$, combien
d'équipes au minimum faut-il prévoir ? Justifier.

\exo{}\diff{2} Un voyageur veut visiter Yaoundé, Douala, Kribi et Limbé reliées par
$Y\!-\!D\,(3\,\text{h})$, $D\!-\!K\,(2\,\text{h})$, $D\!-\!L\,(1\,\text{h})$,
$K\!-\!L\,(2\,\text{h})$. Quel est le trajet le plus court de $Y$ à $L$ ?

% ═══════════════════════════════════════════════════════════════
\section{Sujets type Baccalauréat}

\sujetbac[réseau routier, matrice, parcours eulérien]
On considère cinq villes : Yaoundé (Y), Douala (D), Bafoussam (B), Bertoua (R) et
Ngaoundéré (N). Les liaisons routières directes sont : Y–D, Y–B, Y–R, D–B, B–N, R–N.
\begin{enumerate}[label=\arabic*)]
\item Représenter cette situation par un graphe et donner le degré de chaque ville.
\item Écrire la matrice d'adjacence $M$ (ordre $Y,D,B,R,N$). Le graphe est-il connexe ?
\item Un transporteur souhaite emprunter chaque route exactement une fois. Est-ce possible ?
Si oui, entre quelles villes doit-il commencer et finir ?
\item On admet que le coefficient $(Y,N)$ de $M^2$ vaut $2$. Interpréter ce résultat
en termes d'itinéraire.
\end{enumerate}

\sujetbac[plus court chemin par l'algorithme de Dijkstra]
Un livreur basé à Yaoundé (Y) doit atteindre la ville d'Ebolowa (E). Le réseau routier
(distances en km) est : $Y\!-\!A\,(60)$, $Y\!-\!B\,(30)$, $A\!-\!B\,(20)$, $A\!-\!E\,(40)$,
$B\!-\!C\,(50)$, $C\!-\!E\,(30)$, où $A,B,C$ sont des villes intermédiaires.
\begin{enumerate}[label=\arabic*)]
\item Représenter le réseau par un graphe pondéré.
\item Appliquer l'algorithme de Dijkstra à partir de $Y$ : présenter le tableau des distances
provisoires et des sommets figés.
\item En déduire le plus court chemin de $Y$ à $E$ et sa longueur totale.
\item Le livreur emprunte finalement le trajet $Y\!-\!A\!-\!E$. De combien de km ce trajet
est-il plus long que le plus court chemin ?
\end{enumerate}

\sujetbac[coloration et organisation d'un emploi du temps]
Un lycée doit organiser $6$ activités $a_1,\dots,a_6$. Certaines activités partagent des
élèves et ne peuvent donc avoir lieu simultanément ; les incompatibilités sont :
$a_1a_2$, $a_1a_3$, $a_2a_3$, $a_3a_4$, $a_4a_5$, $a_4a_6$, $a_5a_6$.
\begin{enumerate}[label=\arabic*)]
\item Construire le graphe des incompatibilités et donner le degré de chaque sommet.
\item Montrer que le graphe contient deux triangles. Qu'en déduit-on pour $\chi$ ?
\item Proposer une coloration en utilisant le moins de couleurs possible.
\item En déduire le nombre minimal de créneaux horaires et un emploi du temps possible.
\end{enumerate}

\sujetbac[matrice d'adjacence et dénombrement de chemins]
Un réseau social relie cinq utilisateurs $A,B,C,D,E$. Les amitiés sont : $A\!-\!B$, $A\!-\!C$,
$B\!-\!C$, $B\!-\!D$, $C\!-\!D$, $D\!-\!E$.
\begin{enumerate}[label=\arabic*)]
\item Écrire la matrice d'adjacence $M$ dans l'ordre $A,B,C,D,E$ et donner les degrés.
\item Le réseau est-il connexe ? Justifier.
\item On admet que $M^2=\begin{pmatrix}2&1&1&2&0\\1&3&2&1&1\\1&2&3&1&1\\2&1&1&3&0\\0&1&1&0&1\end{pmatrix}$.
Combien existe-t-il de chaînes de longueur $2$ reliant $A$ à $D$ ? Les expliciter.
\item Existe-t-il une chaîne eulérienne dans ce réseau ? Justifier par le théorème d'Euler.
\end{enumerate}

% ═══════════════════════════════════════════════════════════════
\section{Corrigés}

\corrige{de l'exercice 1}
Somme des degrés $=3+3+2+2+1+1=12=2\,\Card(A)$, donc $\Card(A)=6$ arêtes. Les sommets de
degré impair sont les deux de degré $3$ et les deux de degré $1$ : quatre sommets, nombre
\textbf{pair}, conforme à la propriété.

\corrige{de l'exercice 2}
Tous les sommets de degré $3$ dans un graphe simple d'ordre $4$ : chaque sommet est relié
aux $3$ autres, c'est le graphe complet $K_4$. Il a $\dfrac{4\times 3}{2}=6$ arêtes.

\corrige{de l'exercice 3}
Si quatre sommets sont de degré $4$ dans un graphe \emph{simple} d'ordre $5$, chacun est
relié aux quatre autres ; en particulier au cinquième sommet. Ce cinquième sommet est donc
voisin des quatre autres, son degré vaut $4$ et non $1$ : la liste $4,4,4,4,1$ est
\textbf{impossible}.

\corrige{de l'exercice 4}
Le nombre de sommets de degré impair est toujours \textbf{pair}. Avoir exactement $7$ élèves
(nombre impair) ayant un degré impair est donc \textbf{impossible}.

\corrige{de l'exercice 5}
$K_7$ a $\dfrac{7\times 6}{2}=21$ arêtes et chaque sommet est de degré $7-1=6$.

\corrige{de l'exercice 6}
Chaque paire d'équipes joue une fois : c'est le nombre d'arêtes de $K_8$, soit
$\dfrac{8\times 7}{2}=28$ matches.

\corrige{de l'exercice 7}
Arcs : $A\!\to\!B$, $B\!\to\!C$, $A\!\to\!C$, $C\!\to\!A$. Degrés sortants :
$d^{+}(A)=2$ (vers $B,C$), $d^{+}(B)=1$ ($C$), $d^{+}(C)=1$ ($A$). Degrés entrants :
$d^{-}(A)=1$ ($C$), $d^{-}(B)=1$ ($A$), $d^{-}(C)=2$ ($A,B$). On a $\sum d^{+}=2+1+1=4$ et
$\sum d^{-}=1+1+2=4=\Card(A)$ : l'égalité est bien vérifiée.

\corrige{de l'exercice 8}
Le graphe : $A$ relié à $B,C,D$ ; $B$ relié à $A,D$ ; $C$ relié à $A,D$ ; $D$ relié à
$A,B,C$. Degrés : $d(A)=3$, $d(B)=2$, $d(C)=2$, $d(D)=3$ (sommes des lignes). En partant de
$A$ on atteint $B$, $C$, $D$ : tous les sommets, donc le graphe est \textbf{connexe}.

\corrige{de l'exercice 9}
Le degré d'un sommet est le coefficient diagonal de $M^2$ (nombre de chaînes de longueur $2$
d'un sommet à lui-même = nombre de ses voisins). Ici la diagonale est $3,2,2,3$, donc
$d(A)=3$, $d(B)=2$, $d(C)=2$, $d(D)=3$. Le coefficient $(A,D)$ de $M^2$ vaut $2$ : il y a
\textbf{deux} chaînes de longueur $2$ reliant $A$ à $D$.

\corrige{de l'exercice 10}
Arêtes $A\!-\!B$, $A\!-\!C$, $B\!-\!C$, $C\!-\!D$. Dans l'ordre $A,B,C,D$ :
\[
M=\begin{pmatrix}0&1&1&0\\1&0&1&0\\1&1&0&1\\0&0&1&0\end{pmatrix}.
\]
On a $m_{ij}=m_{ji}$ pour tout couple : $M$ est \textbf{symétrique}. (Sommes des lignes :
$2,2,3,1$, c'est-à-dire $d(A),d(B),d(C),d(D)$, cohérent avec la figure 23.1.)

\corrige{de l'exercice 11}
Les sommets de degré impair sont les deux de degré $5$. Il y en a exactement deux : il
n'existe donc \textbf{pas} de cycle eulérien (qui exigerait tous les degrés pairs), mais il
existe une \textbf{chaîne eulérienne} reliant ces deux sommets de degré $5$ (départ à l'un,
arrivée à l'autre).

\corrige{de l'exercice 12}
Sommets : $A,B,C,D$ (rectangle), $E$ (sommet du triangle au-dessus de $AB$). Arêtes :
$AB,BC,CD,DA$ (rectangle), $AC,BD$ (diagonales), $AE,BE$ (triangle). Degrés :
$d(A)=4$ ($B,D,C,E$), $d(B)=4$ ($A,C,D,E$), $d(C)=3$ ($B,D,A$), $d(D)=3$ ($C,A,B$),
$d(E)=2$ ($A,B$). Deux sommets de degré impair ($C$ et $D$), le graphe est connexe : il
existe une \textbf{chaîne eulérienne} commençant en $C$ et finissant en $D$ (ou l'inverse).
Le dessin se trace donc « sans lever le crayon », mais sans revenir au point de départ.

\corrige{de l'exercice 13}
Degrés $4,4,3,3,2,2$ : exactement deux sommets de degré impair (les deux de degré $3$). Le
graphe étant connexe, il admet une \textbf{chaîne eulérienne} : le facteur peut passer une
seule fois par chaque rue. Comme ce n'est pas un \emph{cycle} (tous les degrés ne sont pas
pairs), il \textbf{ne revient pas} à son point de départ : il commence et finit aux deux
carrefours de degré $3$.

\corrige{de l'exercice 14}
Arêtes : $A\!-\!B\,(2)$, $A\!-\!C\,(5)$, $B\!-\!C\,(1)$, $B\!-\!D\,(7)$, $C\!-\!D\,(3)$.
Dijkstra depuis $A$ :
\[
\begin{array}{c|cccc|l}
\text{ét.} & A & B & C & D & \text{figé}\\\hline
0 & \underline{0} & \infty & \infty & \infty & A\\
1 & & \underline{2}_{(A)} & 5 & \infty & B\\
2 & & & \underline{3}_{(B)} & 9 & C\\
3 & & & & \underline{6}_{(C)} & D
\end{array}
\]
($C$ : via $B$, $2+1=3<5$ ; $D$ : via $C$, $3+3=6<9$.) Plus court chemin
$A\to D$ : $A-B-C-D$ de longueur $\mathbf{6}$.

\corrige{de l'exercice 15}
Arêtes : $P\!-\!Q\,(4)$, $P\!-\!R\,(2)$, $Q\!-\!R\,(1)$, $Q\!-\!S\,(5)$, $R\!-\!S\,(8)$.
Dijkstra depuis $P$ :
\[
\begin{array}{c|cccc|l}
\text{ét.} & P & Q & R & S & \text{figé}\\\hline
0 & \underline{0} & \infty & \infty & \infty & P\\
1 & & 4 & \underline{2}_{(P)} & \infty & R\\
2 & & \underline{3}_{(R)} & & 10 & Q\\
3 & & & & \underline{8}_{(Q)} & S
\end{array}
\]
($Q$ : via $R$, $2+1=3<4$ ; $S$ : via $Q$, $3+5=8<10$ donné par $R$.) Itinéraire le plus
court $P\to S$ : $P-R-Q-S$, distance $\mathbf{8}$ km.

\corrige{de l'exercice 16}
Graphe des incompatibilités : arêtes $E_1E_2$, $E_1E_3$, $E_2E_3$, $E_3E_4$, $E_4E_5$.
Le triangle $E_1E_2E_3$ ($K_3$) impose au moins $3$ couleurs, donc $\chi\ge 3$. Coloration en
$3$ couleurs : $E_1$ = couleur 1, $E_2$ = couleur 2, $E_3$ = couleur 3, $E_4$ = couleur 1
(non adjacent à $E_1$), $E_5$ = couleur 2. Aucune arête ne joint deux sommets de même
couleur, donc $\chi=3$ : il faut au minimum \textbf{3 créneaux horaires}.

\corrige{de l'exercice 17}
$K_5$ : tous les sommets sont adjacents deux à deux, donc deux quelconques doivent recevoir
des couleurs différentes : $\chi(K_5)=5$. Le cycle de longueur $6$ (pair) est biparti : on
alterne deux couleurs autour de l'hexagone, donc $\chi=2$.

\corrige{de l'exercice 18}
\textbf{1)} Sommets $Y,D,B,R,N$ ; arêtes $YD,YB,YR,DB,BN,RN$. Degrés :
$d(Y)=3$ (vers $D,B,R$), $d(D)=2$ ($Y,B$), $d(B)=3$ ($Y,D,N$), $d(R)=2$ ($Y,N$),
$d(N)=2$ ($B,R$). Contrôle : $3+2+3+2+2=12=2\times 6$ arêtes.

\smallskip
\textbf{2)} Dans l'ordre $Y,D,B,R,N$ :
\[
M=\begin{pmatrix}
0&1&1&1&0\\
1&0&1&0&0\\
1&1&0&0&1\\
1&0&0&0&1\\
0&0&1&1&0
\end{pmatrix}.
\]
Elle est symétrique. En partant de $Y$ on atteint $D,B,R$ puis $N$ (via $B$ ou $R$) : tous
les sommets sont accessibles, le graphe est \textbf{connexe}.

\smallskip
\textbf{3)} Parcourir chaque route une seule fois = chercher une chaîne eulérienne. Les
degrés sont $3,2,3,2,2$ : exactement deux sommets de degré impair, $Y$ et $B$. D'après le
théorème d'Euler, une chaîne eulérienne existe et doit \textbf{commencer en $Y$ et finir en
$B$} (ou l'inverse). Un parcours possible : $Y\!-\!R\!-\!N\!-\!B\!-\!D\!-\!Y\!-\!B$.

\smallskip
\textbf{4)} Le coefficient $(Y,N)$ de $M^2$ compte les chaînes de longueur $2$ (deux routes)
reliant Yaoundé à Ngaoundéré. Il vaut $2$ : il y a \textbf{deux itinéraires à deux routes},
qui correspondent aux deux voisins communs de $Y$ et $N$, à savoir $B$ et $R$. Ce sont
$Y\!-\!B\!-\!N$ et $Y\!-\!R\!-\!N$. (Il n'existe en revanche aucune route directe
$Y\!-\!N$.)

\corrige{du sujet 1}
\textbf{1)} Sommets $Y,D,B,R,N$ ; arêtes $YD,YB,YR,DB,BN,RN$. Degrés : $d(Y)=3$, $d(D)=2$,
$d(B)=3$, $d(R)=2$, $d(N)=2$. Somme $12=2\times6$.\\
\textbf{2)} Dans l'ordre $Y,D,B,R,N$, $M=\begin{pmatrix}0&1&1&1&0\\1&0&1&0&0\\1&1&0&0&1\\
1&0&0&0&1\\0&0&1&1&0\end{pmatrix}$, symétrique. En partant de $Y$ on atteint tous les sommets :
le graphe est \textbf{connexe}.\\
\textbf{3)} Deux sommets de degré impair ($Y$ et $B$), graphe connexe : par le théorème
d'Euler, il existe une \textbf{chaîne eulérienne} commençant en $Y$ et finissant en $B$
(ou l'inverse) ; pas de cycle eulérien car tous les degrés ne sont pas pairs.\\
\textbf{4)} $(Y,N)$ de $M^2$ vaut $2$ : il y a deux chaînes de longueur $2$ entre $Y$ et $N$,
à savoir $Y\!-\!B\!-\!N$ et $Y\!-\!R\!-\!N$ (les deux voisins communs $B$ et $R$).

\corrige{du sujet 2}
\textbf{1)} Sommets $Y,A,B,C,E$ ; arêtes pondérées $YA\,(60)$, $YB\,(30)$, $AB\,(20)$,
$AE\,(40)$, $BC\,(50)$, $CE\,(30)$.\\
\textbf{2)} Dijkstra depuis $Y$ :
\[
\begin{array}{c|ccccc|l}
\text{ét.} & Y & A & B & C & E & \text{figé}\\\hline
0 & \underline{0} & \infty & \infty & \infty & \infty & Y\\
1 & & 60 & \underline{30}_{(Y)} & \infty & \infty & B\\
2 & & \underline{50}_{(B)} & & 80 & \infty & A\\
3 & & & & \underline{80}_{(B)} & 90 & C\\
4 & & & & & \underline{90}_{(A)} & E
\end{array}
\]
Détails : $A$ via $B$ : $30+20=50<60$ ; $C$ via $B$ : $30+50=80$ ; $E$ via $A$ : $50+40=90$ ;
via $C$ : $80+30=110>90$, donc $E$ reste à $90$ par $A$.\\
\textbf{3)} En remontant les prédécesseurs : $E\leftarrow A\leftarrow B\leftarrow Y$. Plus
court chemin $Y\to E$ : $Y-B-A-E$ de longueur $30+20+40=\mathbf{90}$ km.\\
\textbf{4)} Le trajet $Y-A-E$ mesure $60+40=100$ km, soit $100-90=\mathbf{10}$ km de plus que
le plus court chemin.

\corrige{du sujet 3}
\textbf{1)} Arêtes $a_1a_2,a_1a_3,a_2a_3,a_3a_4,a_4a_5,a_4a_6,a_5a_6$. Degrés :
$d(a_1)=2$ ($a_2,a_3$), $d(a_2)=2$ ($a_1,a_3$), $d(a_3)=3$ ($a_1,a_2,a_4$),
$d(a_4)=3$ ($a_3,a_5,a_6$), $d(a_5)=2$ ($a_4,a_6$), $d(a_6)=2$ ($a_4,a_5$). Somme
$2+2+3+3+2+2=14=2\times7$ arêtes.\\
\textbf{2)} $a_1a_2a_3$ forment un triangle (arêtes $a_1a_2,a_1a_3,a_2a_3$) et $a_4a_5a_6$
aussi ($a_4a_5,a_4a_6,a_5a_6$). Chaque triangle est un $K_3$, donc $\chi\ge 3$.\\
\textbf{3)} Coloration en $3$ couleurs : $a_1=1$, $a_2=2$, $a_3=3$, $a_4=1$ (adjacent
seulement à $a_3=3$ parmi les déjà colorés), $a_5=2$, $a_6=3$. Vérification : toutes les
arêtes joignent des couleurs distinctes. Donc $\chi=3$ (minoration $\ge3$ et coloration en
$3$).\\
\textbf{4)} Il faut au minimum \textbf{3 créneaux horaires}. Un emploi du temps possible :
créneau 1 $=\{a_1,a_4\}$, créneau 2 $=\{a_2,a_5\}$, créneau 3 $=\{a_3,a_6\}$ (chaque créneau
ne contient que des activités compatibles).

\corrige{du sujet 4}
\textbf{1)} Dans l'ordre $A,B,C,D,E$, arêtes $AB,AC,BC,BD,CD,DE$ :
\[
M=\begin{pmatrix}0&1&1&0&0\\1&0&1&1&0\\1&1&0&1&0\\0&1&1&0&1\\0&0&0&1&0\end{pmatrix}.
\]
Degrés (sommes des lignes) : $d(A)=2$, $d(B)=3$, $d(C)=3$, $d(D)=3$, $d(E)=1$. Contrôle :
$2+3+3+3+1=12=2\times6$ arêtes.\\
\textbf{2)} En partant de $A$ : $A\to B,C$ ; puis $B\to D$ ; puis $D\to E$. Tous les sommets
sont atteints : le réseau est \textbf{connexe}.\\
\textbf{3)} Le coefficient $(A,D)$ de $M^2$ vaut $2$ : il y a \textbf{deux} chaînes de
longueur $2$ de $A$ à $D$. Les voisins communs de $A$ et $D$ sont $B$ et $C$, d'où les
chaînes $A\!-\!B\!-\!D$ et $A\!-\!C\!-\!D$.\\
\textbf{4)} Degrés : $2,3,3,3,1$. Les sommets de degré impair sont $B,C,D,E$ : ils sont au
nombre de \textbf{quatre}. Le théorème d'Euler exige au plus deux sommets de degré impair :
il n'existe donc \textbf{aucune chaîne eulérienne} dans ce réseau.
